\section{Morfogénesis}\label{sec:exp3}

En los experimentos anteriores se conoce la evolución completa y el autómata celular neuronal aprende la transición entre cada par de pasos consecutivos. En este experimento solo se conocen la condición inicial y la configuración final: el autómata debe encontrar una regla local que, aplicada durante cierto número de pasos, lleve de una a la otra.

\TODO{Definir las configuraciones inicial y final a reproducir.}

\TODO{Agregar un modo de entrenamiento en el que la pérdida se calcule sobre la configuración obtenida después de $N$ pasos y no sobre cada transición, ya que el \textit{script} de entrenamiento actual usa pares de pasos consecutivos.}

\TODO{Entrenar, analizar la evolución intermedia que encuentra el autómata y su estabilidad después de alcanzar la configuración final.}
